problem:  
Let \(X\) be a Calabi–Yau threefold with mirror \(X^\vee\).  Let \(\mathcal{F}_A(X^\vee)\) and \(\mathcal{F}_B(X)\) denote the formal germ of the A‑model and B‑model Frobenius manifolds at their respective large limit points, each governed by the genus‑zero WDVV equations and determined up to quadratic changes of coordinates.  

Let \(L_A, L_B \subset \mathcal{H}\) be the corresponding Lagrangian cones in Givental’s symplectic vector space \(\mathcal{H}\) constructed from the A‑ and B‑model genus‑zero theories. Mirror symmetry predicts the existence of a *symplectic transformation* \(T_X:\mathcal{H}\to\mathcal{H}\) such that \(T_X(L_A)=L_B\).  However, \(T_X\) is not known to be canonical outside toric or complete‑intersection cases.  

**Problem:**  
Does there exist, for every Calabi–Yau threefold \(X\), a *canonical symplectic connection* \(\nabla^{\mathrm{mir}}\) on Givental’s symplectic space \(\mathcal{H}\) such that the following holds:

1. The Lagrangian cones \(L_A\) and \(L_B\) are flat (covariantly constant) sections of \(\nabla^{\mathrm{mir}}\) corresponding, respectively, to the A‑ and B‑model Frobenius structures.  
2. Parallel transport along \(\nabla^{\mathrm{mir}}\) from the “A‑model region” to the “B‑model region” induces the mirror map and identifies their Yukawa cubic forms.  
3. The connection \(\nabla^{\mathrm{mir}}\) is unique under the requirement that its curvature form encodes the genus‑zero variation of the BCOV Kodaira–Spencer field theory on the B‑side.

Equivalently, can one find a *global flat symplectic structure* \((\mathcal{H},\nabla^{\mathrm{mir}},\omega)\) underlying both A‑ and B‑model genus‑zero geometries, providing a differential‑geometric realization of mirror symmetry as flat transport between their Frobenius structures?

Why is it a "good" problem:  
This version is mathematically sharper: Givental’s symplectic formalism is well defined for genus‑zero Gromov–Witten and Barannikov–Kontsevich theories, and the proposed object—a flat symplectic connection linking the A‑ and B‑model cones—is a concrete geometric structure whose existence is unknown.  If such a connection exists, it would globalize mirror symmetry at the level of Frobenius geometry, replacing ad hoc identifications by intrinsic parallel transport.  The problem connects enumerative geometry, Hodge theory, and quantization theory in a precise way, inviting tools from symplectic geometry and deformation theory.  It is simple to state yet profound, asking whether the deep equivalence predicted by mirror symmetry can be encoded as curvature‑zero geometry on a single symplectic vector bundle encompassing both models.